\documentclass{article}
\usepackage{amsmath, amssymb, amsthm}
\usepackage{enumitem}
\newtheorem{theorem}{Theorem}[section]
\newtheorem{lemma}[theorem]{Lemma}
\newtheorem{proposition}[theorem]{Proposition}
\newtheorem{corollary}[theorem]{Corollary}
\theoremstyle{definition}
\newtheorem{definition}[theorem]{Definition}
\newtheorem{example}[theorem]{Example}
\newtheorem{remark}[theorem]{Remark}

\begin{document}

\section*{Chapter 2}
\section{Beginnings}

In this chapter, we define quaternion algebras over fields by giving a multiplication table, following Hamilton; we then consider the classical application of understanding rotations in $\mathbb{R}^3$.

\subsection{Conventions}

Throughout this text (unless otherwise stated), we let $F$ be a (commutative) field with algebraic closure $F^{\text{al}}$.

When $G$ is a group, and $H \subseteq G$ is a subset, we write $H \leq G$ when $H$ is a subgroup and $H \trianglelefteq G$ when $H$ is a normal subgroup; if $G$ is abelian (written multiplicatively), we write $G^n := \{g^n : g \in G\} \leq G$ for the subgroup of $n$th powers for $n \in \mathbb{Z}_{>0}$.

We suppose throughout that all rings are associative, not necessarily commutative, with multiplicative identity 1, and that ring homomorphisms preserve 1. In particular, a subring of a ring has the same 1. For a ring $A$, we write $A^\times$ for the multiplicative group of units of $A$. An \textbf{algebra} over the field $F$ is a ring $B$ equipped with a homomorphism $F \to B$ such that the image of $F$ lies in the \textbf{center} $Z(B)$ of $B$, defined by
$$Z(B) := \{\alpha \in B : \alpha\beta = \beta\alpha \text{ for all } \beta \in B\}; \tag{2.1.1}$$
if $Z(B) = F$, we say $B$ is \textbf{central} (as an $F$-algebra). We write $\mathrm{M}_n(F)$ for the $F$-algebra of $n \times n$-matrices with entries in $F$.

One may profitably think of an $F$-algebra as being an $F$-vector space that is also compatibly a ring. If the $F$-algebra $B$ is not the zero ring, then its structure map $F \to B$ is necessarily injective (since 1 maps to 1) and we identify $F$ with its image; keeping track of the structure map just litters notation. The \textbf{dimension} $\dim_F B$ of an $F$-algebra $B$ is its dimension as an $F$-vector space.

A \textbf{homomorphism} of $F$-algebras is a ring homomorphism which restricts to the identity on $F$. An $F$-algebra homomorphism is necessarily $F$-linear. An $F$-algebra homomorphism $B \to B$ is called an \textbf{endomorphism}. By convention (and as usual for functions), endomorphisms act on the left. An invertible $F$-algebra homomorphism $B \xrightarrow{\sim} B'$ is called an \textbf{isomorphism}, and an invertible endomorphism is an \textbf{automorphism}.

The set of automorphisms of $B$ forms a group, which we write as $\operatorname{Aut}(B)$---these maps are necessarily $F$-linear, but we do not include this in the notation. We reserve the notation $\operatorname{End}_F(V)$ for the ring of $F$-linear endomorphisms of the $F$-vector space $V$, and $\operatorname{Aut}_F(V)$ for the group of $F$-linear automorphisms of $V$; in particular, $\operatorname{End}_F(B) \simeq \mathrm{M}_n(F)$ if $n = \dim_F B$.

\begin{remark}
Throughout, whenever we define a homomorphism of objects, we adopt the (categorical) convention extending this to the terms \emph{endomorphism} (homomorphism with equal domain and codomain), \emph{isomorphism} (invertible homomorphism), and \emph{automorphism} (invertible endomorphism).
\end{remark}

A \textbf{division ring} (also called a \textbf{skew field}) is a ring $D$ in which every nonzero element has a (two-sided) inverse, i.e., $D \setminus \{0\}$ is a group under multiplication. A \textbf{division algebra} is an algebra that is a division ring.

\subsection{Quaternion algebras}

In this section, we define quaternion algebras in a direct way, via generators and relations. Throughout the rest of this chapter, suppose that $\operatorname{char} F \neq 2$; the case $\operatorname{char} F = 2$ is treated in Chapter 6.

\begin{definition}
An algebra $B$ over $F$ is a \textbf{quaternion algebra} if there exist $i, j \in B$ such that $1, i, j, ij$ is an $F$-basis for $B$ and
$$i^2 = a, \ j^2 = b, \ \text{and} \ ji = -ij \tag{2.2.2}$$
for some $a, b \in F^\times$.
\end{definition}

The entire multiplication table for a quaternion algebra is determined by the multiplication rules (2.2.2), linearity, and associativity: for example,
$$(ij)^2 = (ij)(ij) = i(ji)j = i(-ij)j = -(i^2)(j^2) = -ab$$
and $j(ij) = (-ij)j = -bi$. Conversely, given $a, b \in F^\times$, one can write down the unique possible associative multiplication table on the basis $1, i, j, k$ compatible with (2.2.2), and then verify independently that it is associative (Exercise 2.4). Accordingly, for $a, b \in F^\times$, we define $\left(\frac{a, b}{F}\right)$ to be the quaternion algebra over $F$ with $F$-basis $1, i, j, ij$ subject to the multiplication (2.2.2); we will also write $(a, b \mid F)$ when convenient for formatting. By definition, we have $\dim_F (a, b \mid F) = 4$.

The map which interchanges $i$ and $j$ gives an isomorphism $\left(\frac{a, b}{F}\right) \simeq \left(\frac{b, a}{F}\right)$, so Definition 2.2.1 is symmetric in $a, b$. The elements $a, b$ are far from unique in determining the isomorphism class of a quaternion algebra: see Exercise 2.3.

If $K \supseteq F$ is a field extension of $F$, then there is a canonical isomorphism
$$\left(\frac{a, b}{F}\right) \otimes_F K \simeq \left(\frac{a, b}{K}\right)$$
extending scalars (same basis, but now spanning a $K$-vector space), so Definition 2.2.1 behaves well with respect to inclusion of fields.

\begin{example}
The $\mathbb{R}$-algebra $\mathbb{H} := \left(\frac{-1, -1}{\mathbb{R}}\right)$ is the ring of quaternions over the real numbers, discovered by Hamilton; we call $\mathbb{H}$ the ring of (real) Hamiltonians (also known as Hamilton's quaternions).
\end{example}

\begin{example}
The ring $\mathrm{M}_2(F)$ of $2 \times 2$-matrices with coefficients in $F$ is a quaternion algebra over $F$: there is an isomorphism $\left(\frac{1, 1}{F}\right) \xrightarrow{\sim} \mathrm{M}_2(F)$ of $F$-algebras induced by
$$i \mapsto \begin{pmatrix} 1 & 0 \\ 0 & -1 \end{pmatrix}, \ j \mapsto \begin{pmatrix} 0 & 1 \\ 1 & 0 \end{pmatrix}.$$
If $F = F^{\mathrm{al}}$ is algebraically closed and $B$ is a quaternion algebra over $F$, then necessarily $B \simeq \mathrm{M}_2(F)$ (Exercise 2.3). Consequently, every quaternion algebra $B$ over $F$ has $B \otimes_F F^{\mathrm{al}} \simeq \mathrm{M}_2(F^{\mathrm{al}})$.
\end{example}

A quaternion algebra $B$ is generated by the elements $i, j$ by definition (2.2.2). However, exhibiting an algebra by generators and relations (instead of by a multiplication table) can be a bit subtle, as the dimension of such an algebra is not a priori clear. But working with presentations is quite useful; and at least for quaternion algebras, we can think in these terms as follows.

\begin{lemma}
An $F$-algebra $B$ is a quaternion algebra if and only if there exist nonzero $i, j \in B$ that generate $B$ as an $F$-algebra and satisfy
$$i^2 = a, \ j^2 = b, \ \text{and} \ ij = -ji \tag{2.2.6}$$
with $a, b \in F^\times$.
\end{lemma}

In other words, once the relations (2.2.6) are satisfied for generators $i, j$, then automatically $B$ has dimension 4 as an $F$-vector space, with $F$-basis $1, i, j, ij$.

\begin{proof}
It is necessary and sufficient to prove that the elements $1, i, j, ij$ are linearly independent. Suppose that $\alpha = t + xi + yj + zij = 0$ with $t, x, y, z \in F$. Using the relations given, we compute that
$$0 = i(\alpha i + i\alpha) = 2a(t + xi).$$
Since $\operatorname{char} F \neq 2$ and $a \neq 0$, we conclude that $t + xi = 0$. Repeating with $j$ and $ij$, we similarly find that $t + yj = t + zij = 0$. Thus
$$\alpha - (t + xi) - (t + yj) - (t + zij) = -2t = 0.$$
Since $i, j$ are nonzero, $B$ is not the zero ring, so $1 \neq 0$; thus $t = 0$ and so $xi = yj = zij = 0$. Finally, if $x \neq 0$, then $i = 0$ so $i^2 = 0 = a$, impossible; hence $x = 0$. Similarly, $y = z = 0$.
\end{proof}

Accordingly, we will call elements $i, j \in B$ satisfying (2.2.6) standard generators for a quaternion algebra $B$.

\begin{remark}
Invertibility of both $a$ and $b$ in $F$ is needed for Lemma 2.2.5: the commutative algebra $B = F[i, j]/(i, j)^2$ is generated by the elements $i, j$ satisfying $i^2 = j^2 = ij = -ji = 0$ but $B$ is \emph{not} a quaternion algebra.
\end{remark}

\begin{remark}
In light of Lemma 2.2.5, we will often drop the symbol $k = ij$ and reserve it for other use. (In particular, in later sections we will want $k$ to represent other quaternion elements.) If we wish to use this abbreviation, we will assign $k := ij$.
\end{remark}

\subsection{Matrix representations}

Every quaternion algebra can be viewed as a subalgebra of $2 \times 2$-matrices over an at most quadratic extension; this is sometimes taken to be the definition!

\begin{proposition}
Let $B := \left(\frac{a, b}{F}\right)$ be a quaternion algebra over $F$ and let $F(\sqrt{a})$ be a splitting field over $F$ for the polynomial $x^2 - a$, with root $\sqrt{a} \in F(\sqrt{a})$. Then the map
\begin{equation}
\begin{aligned}
\lambda \colon B &\to \mathrm{M}_2(F(\sqrt{a})) \\
i, j &\mapsto \begin{pmatrix} \sqrt{a} & 0 \\ 0 & -\sqrt{a} \end{pmatrix}, \begin{pmatrix} 0 & b \\ 1 & 0 \end{pmatrix} \\
t + xi + yj + zij &\mapsto \begin{pmatrix} t + x\sqrt{a} & b(y + z\sqrt{a}) \\ y - z\sqrt{a} & t - x\sqrt{a} \end{pmatrix}
\end{aligned} \tag{2.3.2}
\end{equation}
is an injective $F$-algebra homomorphism and an isomorphism onto its image.
\end{proposition}

\begin{proof}
Injectivity follows by checking $\ker \lambda = \{0\}$ on matrix entries, and the homomorphism property can be verified directly, checking the multiplication table (Exercise 2.6).
\end{proof}

\begin{remark}
Proposition 2.3.1 can be turned around to assert the existence of quaternion algebras: one can check that the set
$$\left\{ \begin{pmatrix} t + x\sqrt{a} & b(y + z\sqrt{a}) \\ y - z\sqrt{a} & t - x\sqrt{a} \end{pmatrix} : t, x, y, z \in F \right\} \subseteq \mathrm{M}_2(F(\sqrt{a}))$$
is an $F$-vector subspace of dimension 4, closed under multiplication, with the matrices $\lambda(i), \lambda(j)$ satisfying the defining relations (2.2.2).
\end{remark}

\textbf{2.3.4.} If $a \notin F^{\times 2}$, then $K = F(\sqrt{a}) \supseteq F$ is a quadratic extension of $F$. Let $\mathrm{Gal}(K \mid F) = \mathrm{Aut}_F(K) \simeq \mathbb{Z}/2\mathbb{Z}$ be the Galois group of $K$ over $F$ and let $\sigma \in \mathrm{Gal}(K \mid F)$ be the nontrivial element. Then we can rewrite the image $\lambda(B)$ in (2.3.2) as
$$\lambda(B) = \left\{ \begin{pmatrix} u & bv \\ \sigma(v) & \sigma(u) \end{pmatrix} : u, v \in K \right\} \subset \mathrm{M}_2(K). \tag{2.3.5}$$

\begin{corollary}
We have an isomorphism
\begin{equation}
\begin{array}{l}
\left(\dfrac{1, b}{F}\right) \xrightarrow{\sim} \mathrm{M}_{2}(F) \\[6pt]
i, j \mapsto \begin{pmatrix} 1 & 0 \\ 0 & -1 \end{pmatrix}, \begin{pmatrix} 0 & b \\ 1 & 0 \end{pmatrix}
\end{array} \tag{2.3.7}
\end{equation}
\end{corollary}

\begin{proof}
Specializing Proposition 2.3.1, we see the map is an injective $F$-algebra homomorphism, so since $\dim_F B = \dim_F \mathrm{M}_2(F) = 4$, the map is also surjective.
\end{proof}

The provenance of the map (2.3.2) is itself important, so we now pursue another (more natural) proof of Proposition 2.3.1.

\textbf{2.3.8.} Let
$$K := F[i] = F \oplus Fi \simeq F[x]/(x^2 - a)$$
be the (commutative) $F$-algebra generated by $i$. Suppose first that $K$ is a field (so $a \notin F^{\times 2}$): then $K \simeq F(\sqrt{a})$ is a quadratic field extension of $F$. The algebra $B$ has the structure of a right $K$-vector space of dimension 2, with basis 1, $j$: explicitly,
$$\alpha = t + xi + yj + zij = (t + xi) + j(y - zi) \in K \oplus jK$$
for all $\alpha \in B$, so $B = K \oplus jK$. We then define the \textbf{left regular representation} of $B$ over $K$ by
\begin{equation}
\begin{array}{l}
\lambda: B \to \operatorname{End}_K(B) \\
\alpha \mapsto (\lambda_\alpha: \beta \mapsto \alpha\beta).
\end{array} \tag{2.3.9}
\end{equation}
Each map $\lambda_\alpha$ is indeed a $K$-linear endomorphism in $B$ (considered as a right $K$-vector space) by associativity in $B$: for all $\alpha, \beta \in B$ and $w \in K$,
$$\lambda_\alpha(\beta w) = \alpha(\beta w) = (\alpha\beta)w = \lambda_\alpha(\beta)w.$$
Similarly, $\lambda$ is an $F$-algebra homomorphism: for all $\alpha, \beta, \nu \in B$
$$\lambda_{\alpha\beta}(\nu) = (\alpha\beta)\nu = \alpha(\beta(\nu)) = (\lambda_\alpha\lambda_\beta)(\nu)$$
reading functions from right to left as usual. The map $\lambda$ is injective ($\lambda$ is a \textbf{faithful} representation) since $\lambda_\alpha = 0$ implies $\lambda_\alpha(1) = \alpha = 0$.

In the basis 1, $j$ we have $\operatorname{End}_K(B) \simeq \mathrm{M}_2(K)$, and $\lambda$ is given by
\begin{equation}
i \mapsto \lambda_i = \begin{pmatrix} i & 0 \\ 0 & -i \end{pmatrix}, \quad j \mapsto \lambda_j = \begin{pmatrix} 0 & b \\ 1 & 0 \end{pmatrix}; \tag{2.3.10}
\end{equation}
these matrices act on column vectors on the left. We then recognize the map $\lambda$ given in (2.3.2).

If $K$ is not a field, then $K \simeq F \times F$, and we repeat the above argument but with $B$ a free module of rank 2 over $K$; then projecting onto one of the factors (choosing $\sqrt{a} \in F$) gives the map $\lambda$, which is still injective and therefore induces an $F$-algebra isomorphism $B \simeq \mathrm{M}_2(F)$.

\begin{remark}
In Proposition 2.3.1, $B$ acts on columns on the left; if instead, one wishes to have $B$ act on the right on rows, give $B$ the structure of a left $K$-vector space and define accordingly the right regular representation instead (taking care about the order of multiplication).
\end{remark}

\textbf{2.3.12.} In some circumstances, it can be notationally convenient to consider variants of the injection (2.3.2): for example
\begin{equation}
\begin{aligned}
B &\rightarrow \mathrm{M}_2(F(\sqrt{a})) \\
t + xi + yj + zij &\mapsto \begin{pmatrix} t + x\sqrt{a} & y + z\sqrt{a} \\ b(y - z\sqrt{a}) & t - x\sqrt{a} \end{pmatrix}
\end{aligned} \tag{2.3.13}
\end{equation}
is obtained by taking the basis $1, b^{-1}j$, equivalently postcomposing by $\begin{pmatrix} 1 & 0 \\ 0 & b \end{pmatrix}$. See also Exercise 2.12.

\begin{remark}
The left regular representation 2.3.8 is not the only way to embed $B$ as a subalgebra of $2 \times 2$-matrices. Indeed, the ``splitting'' of quaternion algebras in this way, in particular the question of whether or not $B \simeq \mathrm{M}_2(F)$, is a theme that will reappear throughout this text. For a preview, see Main Theorem 5.4.4.
\end{remark}

\textbf{2.3.15.} Thinking of a quaternion algebra as in 2.3.8 as a right $K$-vector space suggests notation for quaternion algebras that is also useful: for a peek, see 6.1.5.

\subsection{Rotations}

To conclude this chapter, we return to Hamilton's original design: quaternions model rotations in 3-dimensional space. This development is not only historically important but it also previews many aspects of the general theory of quaternion algebras over fields. In this section, we follow Hamilton and take $k := ij$.

Proposition 2.3.1 provides an $\mathbb{R}$-algebra embedding
\begin{equation}
\begin{aligned}
\lambda : \mathbb{H} &\hookrightarrow \operatorname{End}_{\mathbb{C}}(\mathbb{H}) \simeq \mathrm{M}_2(\mathbb{C}) \\
t + xi + yj + zk &= u + j\bar{v} \mapsto \begin{pmatrix} t + xi & -y - zi \\ y - zi & t - xi \end{pmatrix} = \begin{pmatrix} u & -v \\ \bar{v} & \bar{u} \end{pmatrix}
\end{aligned} \tag{2.4.1}
\end{equation}
where $u := t + xi$ and $v := y + zi$ and $\bar{v}$ denotes complex conjugation. (The abuse of notation, taking $i \in \mathbb{H}$ as well as $i \in \mathbb{C}$ is harmless: we may think of $\mathbb{C} \subset \mathbb{H}$.) We have
$$\det \begin{pmatrix} u & -v \\ \bar{v} & \bar{u} \end{pmatrix} = |u|^2 + |v|^2 = t^2 + x^2 + y^2 + z^2,$$
thus $\mathbb{H}^\times = \mathbb{H} \setminus \{0\}$. If preferred, see (2.3.13) to obtain matrices of the form $\begin{pmatrix} u & v \\ -\bar{v} & \bar{u} \end{pmatrix}$ instead.

\textbf{2.4.2.} We define the subgroup of unit Hamiltonians as
$$\mathbb{H}^1 := \{t + xi + yj + zk \in \mathbb{H} : t^2 + x^2 + y^2 + z^2 = 1\}.$$
(In some contexts, one also writes $\mathrm{GL}_1(\mathbb{H}) = \mathbb{H}^\times$ and $\mathrm{SL}_1(\mathbb{H}) = \mathbb{H}^1$.)

As a set, the unit Hamiltonians are naturally identified with the 3-sphere in $\mathbb{R}^4$. As groups, we have an isomorphism $\mathbb{H}^1 \simeq \mathrm{SU}(2)$ with the special unitary group of rank 2, where
\begin{equation}
\mathrm{SU}(n) := \{A \in \mathrm{SL}_n(\mathbb{C}) : A^* = A^{-1}\} \tag{2.4.3}
\end{equation}
where $A^* = \overline{A}^{\mathrm{t}}$ is the (complex) conjugate transpose of $A$.

\begin{definition}
Let $\alpha \in \mathbb{H}$. We say $\alpha$ is \textbf{real} if $\alpha \in \mathbb{R}$, and we say $\alpha$ is \textbf{pure} (or \textbf{imaginary}) if $\alpha \in \mathbb{R}i + \mathbb{R}j + \mathbb{R}k$.
\end{definition}

\textbf{2.4.5.} Just as for the complex numbers, every element of $\mathbb{H}$ is the sum of its real part and its pure (imaginary) part. And just like complex conjugation, we define a (quaternion) conjugation map
$$\overline{\phantom{\alpha}} : \mathbb{H} \to \mathbb{H}$$
\begin{equation}
\alpha = t + (xi + yj + zk) \mapsto \overline{\alpha} = t - (xi + yj + zk) \tag{2.4.6}
\end{equation}
by negating the imaginary part. We compute that
\begin{equation}
\begin{aligned}
\alpha + \overline{\alpha} &= \mathrm{tr}(\lambda(\alpha)) = 2t \\
\|\alpha\|^2 &:= \det(\lambda(\alpha)) = \alpha\overline{\alpha} = \overline{\alpha}\alpha = t^2 + x^2 + y^2 + z^2.
\end{aligned} \tag{2.4.7}
\end{equation}
The notation $\|\cdot\|^2$ is used to indicate that it agrees with the usual square norm on $\mathbb{H} \simeq \mathbb{R}^4$.

The conjugate transpose map on $\mathrm{M}_2(\mathbb{C})$ restricts to quaternion conjugation on the image of $\mathbb{H}$ in (2.4.1), also known as adjugation
$$\alpha = \begin{pmatrix} u & -v \\ \overline{v} & \overline{u} \end{pmatrix} \mapsto \lambda(\overline{\alpha}) = \begin{pmatrix} \overline{u} & v \\ -\overline{v} & u \end{pmatrix}.$$
Thus the elements $\alpha \in \mathbb{H}$ such that $\lambda(\overline{\alpha}) = \lambda(\alpha)$ (i.e., $A^* = A$, and we say $A$ is Hermitian), are exactly the scalar (real) matrices; and those that are skew-Hermitian, i.e., $A^* = -A$, are exactly the pure quaternions. The conjugation map plays a crucial role for quaternion algebras and is the subject of the next chapter (Chapter 3), where to avoid confusion with other notions of conjugation we refer to it as the standard involution.

\textbf{2.4.8.} Let
$$\mathbb{H}^0 := \{v = xi + yj + zk \in \mathbb{H} : x, y, z \in \mathbb{R}\} \simeq \mathbb{R}^3$$
be the set of pure Hamiltonians, the three-dimensional real space on which we will soon see that the (unit) Hamiltonians act by rotations. (The reader should not confuse $v \in \mathbb{H}^0$ with $v$ the entry of a $2 \times 2$-matrix in a local instantiation above.) For $v \in \mathbb{H}^0 \simeq \mathbb{R}^3$,
\begin{equation}
\|v\|^2 = x^2 + y^2 + z^2 = \det(\lambda(v)), \tag{2.4.9}
\end{equation}
and from (2.4.1),
$$\mathbb{H}^0 = \{v \in \mathbb{H} : \mathrm{tr}(\lambda(v)) = v + \bar{v} = 0\}.$$
We again see that $\bar{v} = -v$ for $v \in \mathbb{H}^0$.

The set $\mathbb{H}^0$ is not closed under multiplication: if $v, w \in \mathbb{H}^0$, then
\begin{equation}
vw = -v \cdot w + v \times w \tag{2.4.10}
\end{equation}
where $v \cdot w$ is the dot product on $\mathbb{R}^3$ and $v \times w \in \mathbb{H}^0$ is the \textbf{cross product}, defined as the determinant
\begin{equation}
v \times w = \det \begin{pmatrix} i & j & k \\ v_1 & v_2 & v_3 \\ w_1 & w_2 & w_3 \end{pmatrix} \tag{2.4.11}
\end{equation}
where $v = v_1 i + v_2 j + v_3 k$ and $w = w_1 i + w_2 j + w_3 k$, so
$$v \cdot w = v_1 w_1 + v_2 w_2 + v_3 w_3$$
and
$$v \times w = (v_2 w_3 - v_3 w_2)i + (v_3 w_1 - v_1 w_3)j + (v_1 w_2 - v_2 w_1)k.$$
The formula (2.4.10) is striking: it contains three different kinds of `multiplications'!

\begin{lemma}
For all $v, w \in \mathbb{H}^0$, the following statements hold.
\begin{enumerate}[label=(\alph*)]
\item $vw \in \mathbb{H}^0$ if and only if $v, w$ are orthogonal.
\item $v^2 = -\|v\|^2 \in \mathbb{R}$.
\item $wv = -vw$ if and only if $v, w$ are orthogonal.
\end{enumerate}
\end{lemma}

\begin{proof}
Apply (2.4.10).
\end{proof}

\textbf{2.4.13.} The group $\mathbb{H}^1$ acts on our three-dimensional space $\mathbb{H}^0$ (on the left) by conjugation:
\begin{equation}
\begin{aligned}
\mathbb{H}^1 \times \mathbb{H}^0 &\to \mathbb{H}^0 \\
v &\mapsto \alpha v \alpha^{-1};
\end{aligned} \tag{2.4.14}
\end{equation}
indeed, $\mathrm{tr}(\lambda(\alpha v \alpha^{-1})) = \mathrm{tr}(\lambda(v)) = 0$ by properties of the trace, so $\alpha v \alpha^{-1} \in \mathbb{H}^0$. Or
$$\mathbb{H}^0 = \{v \in \mathbb{H} : v^2 \in \mathbb{R}_{\leq 0}\}$$
and this latter set is visibly stable under conjugation. The representation (2.4.14) is called the \textbf{adjoint representation}.

\textbf{2.4.15.} Let $\alpha \in \mathbb{H}^1 \setminus \{\pm 1\}$. Then there exists a unique $\theta \in (0, \pi)$ such that
\begin{equation}
\alpha = t + xi + yj + zk = \cos \theta + (\sin \theta)I(\alpha) \tag{2.4.16}
\end{equation}
where $I(\alpha)$ is pure and $\|I(\alpha)\| = 1$: to be precise, we take $\theta$ such that $\cos \theta = t$ and
$$I(\alpha) := \frac{xi + yj + zk}{|\sin \theta|}.$$
We call $I(\alpha)$ the \textbf{axis} of $\alpha$, and observe that $I(\alpha)^2 = -1$.

\begin{remark}
In analogy with Euler's formula, we can write (2.4.16) as
$$\alpha = \exp(I(\alpha)\theta).$$
\end{remark}

We are now prepared to identify this action by quaternions with rotations. As usual, let
$$\mathrm{O}(n) := \{A \in \mathrm{M}_n(\mathbb{R}) : A^{\mathrm{t}} = A^{-1}\}$$
be the orthogonal group of $\mathbb{R}^n$ (preserving the standard inner product), and let
$$\mathrm{SO}(n) := \{A \in \mathrm{O}(n) : \det(A) = 1\} \trianglelefteq \mathrm{O}(n)$$
be the special orthogonal group of rotations of $\mathbb{R}^n$, a normal subgroup of index 2 fitting into the exact sequence
$$1 \to \mathrm{SO}(n) \to \mathrm{O}(n) \xrightarrow{\det} \{\pm 1\} \to 1.$$

\begin{proposition}
$\mathbb{H}^1$ acts by rotation on $\mathbb{H}^0 \simeq \mathbb{R}^3$ via conjugation (2.4.14): specifically, $\alpha$ acts by rotation through the angle $2\theta$ about the axis $I(\alpha)$.
\end{proposition}

\begin{proof}
Let $\alpha \in \mathbb{H}^1 \setminus \{\pm 1\}$. Then for all $v \in \mathbb{H}^0$,
$$\|\alpha v \alpha^{-1}\|^2 = \|v\|^2$$
by (2.4.9), so $\alpha$ acts by a matrix belonging to $\mathrm{O}(3)$.

But we can be more precise. Let $j' \in \mathbb{H}^0$ be a unit vector orthogonal to $i' = I(\alpha)$. Then $(i')^2 = (j')^2 = -1$ by Lemma 2.4.12(b) and $j'i' = -i'j'$ by Lemma 2.4.12(c), so without loss of generality we may suppose that $I(\alpha) = i$ and $j' = j$. Thus $\alpha = t + xi = \cos\theta + (\sin\theta)i$ with $t^2 + x^2 = \cos^2\theta + \sin^2\theta = 1$, and $\alpha^{-1} = t - xi$.

We have $\alpha i \alpha^{-1} = i$, and
\begin{equation}
\begin{aligned}
\alpha j \alpha^{-1} &= (t + xi)j(t - xi) = (t + xi)(t + xi)j \\
&= ((t^2 - x^2) + 2txi)j = (\cos 2\theta)j + (\sin 2\theta)k
\end{aligned} \tag{2.4.19}
\end{equation}
by the double angle formula. Consequently,
$$\alpha k \alpha^{-1} = i(\alpha j \alpha^{-1}) = (-\sin 2\theta)j + (\cos 2\theta)k$$
so the matrix of $\alpha$ in the basis $i, j, k$ is
\begin{equation}
A = \begin{pmatrix}
1 & 0 & 0 \\
0 & \cos 2\theta & -\sin 2\theta \\
0 & \sin 2\theta & \cos 2\theta
\end{pmatrix}, \tag{2.4.20}
\end{equation}
a (counterclockwise) rotation (determinant 1) through the angle $2\theta$ about $i$.
\end{proof}

\begin{corollary}
The action (2.4.13) defines a group homomorphism $\mathbb{H}^1 \to \mathrm{SO}(3)$, fitting into an exact sequence
$$1 \to \{\pm 1\} \to \mathbb{H}^1 \to \mathrm{SO}(3) \to 1.$$
\end{corollary}

\begin{proof}
The map $\mathbb{H}^1 \rightarrow \mathrm{SO}(3)$ is surjective, since every element of $\mathrm{SO}(3)$ is rotation about some axis (Exercise 2.15). If $\alpha$ belongs to the kernel, then $\alpha = \cos \theta + (\sin \theta)I(\alpha)$ must have $\sin \theta = 0$ so $\alpha = \pm 1$.
\end{proof}

\textbf{2.4.22.} The matrix representation of $\mathbb{H}$ in section 2.4 extends to a matrix representation of $\mathbb{H} \otimes_{\mathbb{R}} \mathbb{C}$, and this representation and its connection to unitary matrices is still used widely in quantum mechanics. In the embedding with
$$i \mapsto \begin{pmatrix} i & 0 \\ 0 & -i \end{pmatrix}, \quad -j \mapsto \begin{pmatrix} 0 & 1 \\ -1 & 0 \end{pmatrix}, \quad -k \mapsto \begin{pmatrix} 0 & i \\ i & 0 \end{pmatrix}$$
whose images are unitary matrices, we multiply by $-i$ to obtain Hermitian matrices
$$\sigma_z := \begin{pmatrix} 1 & 0 \\ 0 & -1 \end{pmatrix}, \quad \sigma_y := \begin{pmatrix} 0 & -i \\ i & 0 \end{pmatrix}, \quad \sigma_x := \begin{pmatrix} 0 & 1 \\ 1 & 0 \end{pmatrix}$$
where $\sigma_x, \sigma_y, \sigma_z$ are the famous \textbf{Pauli spin matrices}. Because of this application to the spin (a kind of angular momentum) of an electron in particle physics, the group $\mathbb{H}^1$ also goes by the name $\mathbb{H}^1 \simeq \mathrm{Spin}(3)$.

The extra bit of information conveyed by spin can also be seen by the ``belt trick'' [Hans2006, Chapter 2].

\textbf{2.4.23.} We conclude with one final observation, returning to the formula (2.4.10). There is another way to mix the dot product and cross product (2.4.11) in $\mathbb{H}$: we define the \textbf{scalar triple product}
\begin{equation}
\begin{aligned}
\mathbb{H} \times \mathbb{H} \times \mathbb{H} &\rightarrow \mathbb{R} \\
(u, v, w) &\mapsto u \cdot (v \times w).
\end{aligned} \tag{2.4.24}
\end{equation}
Amusingly, this gives a way to ``multiply'' \emph{triples} of triples! The map (2.4.24) defines an alternating, trilinear form (Exercise 2.19). If $u, v, w \in \mathbb{H}^0$, then the scalar triple product is a determinant
$$u \cdot (v \times w) = \det \begin{pmatrix} u_1 & u_2 & u_3 \\ v_1 & v_2 & v_3 \\ w_1 & w_2 & w_3 \end{pmatrix}$$
and $|u \cdot (v \times w)|$ is the volume of a parallelepiped in $\mathbb{R}^3$ whose sides are given by $u, v, w$.

\subsection*{Exercises}

Let $F$ be a field with $\operatorname{char} F \neq 2$.

\begin{enumerate}
\item Show that a (not necessarily associative) $F$-algebra is associative if and only if the associative law holds on a basis, and then check that the multiplication table implied by (2.2.2) is associative.

\item Show that if $B$ is an $F$-algebra generated by $i, j \in B$ and $1, i, j$ are linearly dependent, then $B$ is commutative.

\item Verify directly that the map $\left(\frac{1,1}{F}\right) \xrightarrow{\sim} \mathrm{M}_2(F)$ in Example 2.2.4 is an isomorphism of $F$-algebras.

\item Let $a, b \in F^\times$.
\begin{enumerate}[label=(\alph*)]
\item Show that $\left(\frac{a, b}{F}\right) \simeq \left(\frac{a, -ab}{F}\right) \simeq \left(\frac{b, -ab}{F}\right)$.
\item Show that if $c, d \in F^\times$ then $\left(\frac{a, b}{F}\right) \simeq \left(\frac{ac^2, bd^2}{F}\right)$. Conclude that if $F^\times / F^{\times 2}$ is finite, then there are only finitely many isomorphism classes of quaternion algebras over $F$, and in particular that if $F^{\times 2} = F^\times$ then there is only one isomorphism class $\left(\frac{1,1}{F}\right) \simeq \mathrm{M}_2(F)$. [The converse is not true, see Exercise 3.15.]
\item Show that if $B = \left(\frac{a, b}{\mathbb{R}}\right)$ is a quaternion algebra over $\mathbb{R}$, then $B \simeq \mathrm{M}_2(\mathbb{R})$ or $B \simeq \mathbb{H}$, the latter occurring if and only if $a < 0$ and $b < 0$. Conclude that if $B$ is a division quaternion algebra over $\mathbb{R}$, then $B \simeq \mathbb{H}$.
\item Let $B$ be a quaternion algebra over $F$. Show that $B \otimes_F F^{\mathrm{al}} \simeq \mathrm{M}_2(F^{\mathrm{al}})$, where $F^{\mathrm{al}}$ is an algebraic closure of $F$.
\item Refine part (d) as follows. A field $K \supseteq F$ is a splitting field for $B$ if $B \otimes_F K \simeq \mathrm{M}_2(K)$. Show that $B$ has a splitting field $K$ with $[K : F] \leq 2$.
\end{enumerate}

\item Let $B = \left(\frac{a, b}{F}\right)$ be a quaternion algebra over $F$. Let $i' \in B \setminus F$ satisfy $(i')^2 = a' \in F^\times$. Show that there exists $b' \in F^\times$ and an isomorphism $B \simeq \left(\frac{a', b'}{F}\right)$ (under which $i'$ maps to the first standard generator).

\item Use the quaternion algebra $B = \left(\frac{-1, -1}{F}\right)$, multiplicativity of the determinant, and the left regular representation (2.3.2) to show that if two elements of $F$ can be written as the sum of four squares, then so too can their product (a discovery of Euler in 1748). [In Chapter 3, this statement will follow immediately from the multiplicativity of the reduced norm on $B$; here, the formula is derived easily from multiplicativity of the determinant.]

\item Let $B$ be an $F$-algebra. Show that if $B$ is a quaternion algebra over $F$, then $B$ is central.

\item Let $A, B$ be $F$-algebras, and let $\phi : A \to B$ be a surjective $F$-algebra homomorphism. Show that $\phi$ restricts to an $F$-algebra homomorphism $Z(A) \to Z(B)$.

\item Prove the following partial generalization of Exercise 2.3(b). Let $B$ be a finite-dimensional algebra over $F$.
\begin{enumerate}[label=(\alph*)]
\item Show that every element $\alpha \in B$ satisfies a unique monic polynomial of smallest degree with coefficients in $F$.
\item Suppose that $B = D$ is a division algebra. Show that the minimal polynomial of $\alpha \in D$ is irreducible over $F$. Conclude that if $F = F^{\text{al}}$ is algebraically closed, then $D = F$.
\end{enumerate}

\item Prove Proposition 2.3.1: show directly that the map
$$\lambda: B \rightarrow \mathrm{M}_2(F(\sqrt{a})), \quad i, j \mapsto \begin{pmatrix} \sqrt{a} & 0 \\ 0 & -\sqrt{a} \end{pmatrix}, \begin{pmatrix} 0 & b \\ 1 & 0 \end{pmatrix}$$
extends uniquely to an injective $F$-algebra homomorphism. [Hint: check that the relations are satisfied.]

\item \begin{enumerate}[label=(\alph*)]
\item Show explicitly that every quaternion algebra $B = (a, b \mid F)$ is isomorphic to an $F$-subalgebra of $\mathrm{M}_4(F)$ via the left (or right) regular representation over $F$: write down $4 \times 4$-matrices representing $i$ and $j$ and verify that the relations $i^2 = a, j^2 = b, ji = -ij$ hold for these matrices. Note the $2 \times 2$-block structure of these matrices.
\item With respect to a suitable such embedding in (a) for $B = \mathbb{H}$, verify that the quaternionic conjugation map $\alpha \mapsto \overline{\alpha}$ is the matrix transpose, and the matrix determinant is the \emph{square} of the norm $\|\alpha\|^2 = \alpha \overline{\alpha}$.
\end{enumerate}

\item In certain circumstances, one may not want to ``play favorites'' in the left regular representation (Proposition 2.3.1) and so involve $i$ and $j$ on more equal footing. To this end, show that the map
\begin{equation}
B = \left(\frac{a, b}{F}\right) \rightarrow \mathrm{M}_2(F(\sqrt{a}, \sqrt{b})), \quad
i, j \mapsto \begin{pmatrix} \sqrt{a} & 0 \\ 0 & -\sqrt{a} \end{pmatrix}, \begin{pmatrix} 0 & \sqrt{b} \\ \sqrt{b} & 0 \end{pmatrix} \tag{2.4.25}
\end{equation}
is an injective $F$-algebra isomorphism. How is it related to the left regular representation?

\item Let $B = (a, b \mid F)$ be a quaternion algebra over $F$. For a nonzero element $\alpha = t + xi + yj + zk \in B$, show that the following are equivalent:
\begin{enumerate}[label=(\roman*)]
\item $t = 0$; and
\item $\alpha \notin F$ and $\alpha^2 \in F$.
\end{enumerate}
[So the notion of ``pure quaternion'' is not tethered to a particular basis.]

\item Verify that (2.3.7) is an isomorphism of $F$-algebras, and interpret this map as arising from the left regular representation via a map $B \hookrightarrow \mathrm{M}_2(F \times F) \rightarrow \mathrm{M}_2(F)$.

\item Show that every rotation $A \in \mathrm{SO}(3)$ fixes an axis. [Hint: Consider the eigenvalues of $A$.]

\item For $\nu \in \mathbb{H}^0$ and $\beta \in \mathbb{H}^0 \setminus \{0\}$, consider the map $\nu \mapsto \beta^{-1} \overline{\nu} \beta = -\beta^{-1} \nu \beta \in \mathbb{H}^0$. Show that this map is the reflection across the plane $\{w \in \mathbb{H}^0 : \mathrm{tr}(\lambda(\beta w)) = 0\}$. (For example, taking $\beta = i$, the map is $xi + yj + zk \mapsto -xi + yj + zk$.)

\item In Corollary 2.4.21, we showed that $\mathrm{SU}(2) \simeq \mathbb{H}^1$ has a 2-to-1 map to $\mathrm{SO}(3)$, where $\mathbb{H}^1$ acts on $\mathbb{H}^0 \simeq \mathbb{R}^3$ by conjugation: quaternions model rotations in three-dimensional space, with spin. Quaternions also model rotations in four-dimensional space, as follows.
\begin{enumerate}[label=(\alph*)]
\item Show that the map
\begin{equation}
\begin{array}{c}
(\mathbb{H}^1 \times \mathbb{H}^1) \times \mathbb{H} \to \mathbb{H} \\
x \mapsto \alpha x \beta^{-1}
\end{array} \tag{2.4.26}
\end{equation}
defines a (left) action of $\mathbb{H}^1 \times \mathbb{H}^1$ on $\mathbb{H} \simeq \mathbb{R}^4$, giving a group homomorphism
$$\phi: \mathbb{H}^1 \times \mathbb{H}^1 \to \mathrm{O}(4).$$
\item Show that $\phi$ surjects onto $\mathrm{SO}(4) < \mathrm{O}(4)$. [Hint: If $A \in \mathrm{SO}(4)$ fixes $1 \in \mathbb{H}$, then $A$ restricted to $\mathbb{H}^0$ is a rotation and so is given by conjugation. More generally, if $A1 = \alpha$, consider $x \mapsto \alpha^{-1}Ax$.]
\item Show that the kernel of $\phi$ is $\{\pm 1\}$ embedded diagonally, so there is an exact sequence
$$1 \to \{\pm 1\} \to \mathrm{SU}(2) \times \mathrm{SU}(2) \to \mathrm{SO}(4) \to 1.$$
[More generally, the universal cover of $\mathrm{SO}(n)$ for $n \geq 3$ is a double cover called the spin group $\mathrm{Spin}(n)$, and so Corollary 2.4.21 shows that $\mathrm{Spin}(3) \simeq \mathrm{SU}(2)$ and this exercise shows that $\mathrm{Spin}(4) \simeq \mathrm{SU}(2) \times \mathrm{SU}(2)$. For further reading, see e.g.\ Fulton--Harris [FH91, Lecture 20].]
\end{enumerate}

\item Let $\rho_{u,\theta}: \mathbb{R}^3 \to \mathbb{R}^3$ be the counterclockwise rotation by the angle $\theta$ about the axis $u \in \mathbb{R}^3 \simeq \mathbb{H}^0$, with $\|u\| = 1$. Prove Rodrigues's rotation formula: for all $v \in \mathbb{R}^3$,
$$\rho_{u,\theta}(v) = (\cos \theta)v + (\sin \theta)(u \times v) + (1 - \cos \theta)(u \cdot v)u$$
where $u \times v$ and $u \cdot v$ are the cross and dot product, respectively.

\item Verify that the map (2.4.24) is a trilinear alternating form on $\mathbb{H}$, i.e., show the form is linear when any two of the three arguments are fixed and zero when two arguments are equal.

\item Let $B$ be a quaternion algebra over $F$ and let $\mathrm{M}_2(B)$ be the ring of $2 \times 2$-matrices over $B$. (Be careful in the definition of matrix multiplication: $B$ is noncommutative!) Consider the Cayley determinant:
$$\mathrm{Cdet}: \mathrm{M}_2(B) \to B, \quad \mathrm{Cdet} \begin{pmatrix} \alpha & \beta \\ \gamma & \delta \end{pmatrix} = \alpha \delta - \gamma \beta$$
\begin{enumerate}[label=(\alph*)]
\item Show that $\mathrm{Cdet}$ is $F$-multilinear in the rows and columns of the matrix.
\item Show that $\mathrm{Cdet}$ is not left $B$-multilinear in the rows of the matrix.
\item Give an example showing that $\mathrm{Cdet}$ is not multiplicative.
\item Find a matrix $A \in \mathrm{M}_2(\mathbb{H})$ that is invertible (i.e., having a two-sided inverse) but has $\mathrm{Cdet}(A) = 0$. Then find such an $A$ with the further property that its transpose has nonzero determinant but is not invertible.
\end{enumerate}
[Moral: be careful with matrix rings over noncommutative rings! For more on quaternionic determinants, including the Dieudonn\'{e} determinant, see Aslaksen [Asl96].]
\end{enumerate}

\end{document}